%% series_4_frobenius.tex -- version ampliada de series_4 (los tres casos del metodo de Frobenius)
%% Agregados: clasificacion por la diferencia r1-r2 (recta de raices y diagrama de decision),
%% tabla resumen, graficos de las soluciones con logaritmo, ejemplo sin logaritmo aunque las
%% raices difieran en un entero (Bessel de orden 1/2), tabla de A en funcion de gamma, ejercicios.
%% Correcciones respecto del original:
%%  - la numeracion "Caso 1/2/3" del teorema (real distinto, igual, complejo) chocaba con la del
%%    resumen (no entero, igual, entero); ahora los casos son I, II, III segun r1-r2;
%%  - la condicion de consistencia usa q_{N-k} (el original decia q_{N+k});
%%  - en el Ejemplo con parametro gamma, q(x)=-1/4+(gamma/4)x, de modo que la recurrencia lleva
%%    gamma/4 y el valor crtico es gamma=4 (el original afirmaba gamma=1);
%%  - en el resumen, el segundo caso pide b_0 libre (no "a_0 distinto de 0" para y_2).
%% Requiere en el preambulo del documento principal:
%%   \usepackage{pgfplots}  \pgfplotsset{compat=1.18}
%% Convenciones de notacion del capitulo: en la forma estandar y''+P(x)y'+Q(x)y=0
%% los coeficientes se llaman P y Q; en la forma de Frobenius
%% x^2 y''+x p(x) y'+q(x) y=0 se usan p=(x-x0)P y q=(x-x0)^2 Q (minusculas).
 
%% --- Correccion de \ref a cajas numeradas (ver sistemas_2.tex) --------------
\makeatletter
\ifdefined\bbereffixdone\else
  \def\bbereffixdone{}
  \tcbset{before upper={%
    \ifx\thetcbcounter\@empty\else
      \protected@edef\@currentlabel{\thetcbcounter}%
      \ifdefined\phantomsection\phantomsection\fi
    \fi}}
\fi
\makeatother
 
%% --- Estilo comun de los graficos de este capitulo --------------------------
\pgfplotsset{
  bbeserie/.style={width=.94\linewidth, height=6.2cm, axis lines=middle,
    xlabel=$x$, ylabel=$y$, tick label style={font=\small}, label style={font=\small},
    legend style={font=\footnotesize, draw=bbe!40, fill=white, fill opacity=.9, text opacity=1},
    grid=major, grid style={draw=bbe!12}, clip=true,
    every axis plot/.append style={thick}},
  bbeosc/.style={bbeserie, axis lines=left}
}
\providecommand{\bbezero}[2]{\addplot[bbegray!70, thin, forget plot] coordinates {(#1,0) (#2,0)};}
\tikzset{
  bbecaja/.style={draw=bbe, fill=bbe!8, rounded corners=2pt, text=black, align=center,
    font=\small\sffamily, inner sep=4pt, text width=3.6cm, minimum height=1.1cm},
  bbeflujo/.style={-{Stealth[length=5pt]}, thick}
}
 
 
En la sección anterior vimos cómo se pueden obtener soluciones en serie de Frobenius para la ecuación diferencial
\begin{equation}
x^2 y'' + xp(x)y' + q(x)y = 0, \qquad x > 0.
\label{eq:s4-frobenius}
\end{equation}
 
En esta sección damos alguna justificación para el Teorema de Frobenius y lo extendemos al caso cuando las raíces de la ecuación indicial para \eqref{eq:s4-frobenius} difieren por un número entero. Primero asumiremos que $x > 0$, ya que nuestros resultados pueden extenderse fácilmente a $x < 0$ (reemplazando $x^r$ por $|x|^r$). Comenzamos estableciendo la existencia de al menos una solución en serie de Frobenius.
 
\begin{remark}
    Asumiendo que $x = 0$ es un punto singular regular de \eqref{eq:s4-frobenius}, se deduce que $p$ y $q$ son analíticas en $x = 0$, por lo tanto podemos escribir
\[ p(x) = \sum_{n=0}^{\infty} p_n x^n, \qquad q(x) = \sum_{n=0}^{\infty} q_n x^n, \]
para $|x| < R$. En consecuencia, \eqref{eq:s4-frobenius} se puede escribir como
\[ x^2 y'' + x \Bigl(\sum_{n=0}^{\infty} p_n x^n\Bigr) y' + \Bigl(\sum_{n=0}^{\infty} q_n x^n\Bigr) y = 0. \]
 
Intentamos una solución en serie de Frobenius y, por lo tanto, sea
\[ y(x) = x^r \sum_{n=0}^{\infty} a_n x^n, \qquad a_0 \neq 0, \]
donde $r$ y $a_n$ son constantes por determinar. Sabemos que la ecuación indicial emerge de manera natural,
\begin{equation}
F(r):=r(r - 1) + p_0 r + q_0 = 0,
\label{eq:s4-indicial}
\end{equation}
y que los coeficientes se obtienen de la recurrencia general
\begin{equation}
F(r+n)\,a_n=-\sum_{k=0}^{n-1}\bigl[(r+k)\,p_{n-k}+q_{n-k}\bigr]a_k,\qquad n=1,2,\dots
\label{eq:s4-recgen}
\end{equation}
(la fórmula \eqref{eq:s3-recgen} de la sección anterior). El único obstáculo es que el coeficiente $F(r+n)$ de $a_n$ \emph{se anule}. Como $F$ es un polinomio cuadrático con ceros $r_1$ y $r_2$, eso solo puede ocurrir si $r+n$ es una de las raíces.
\end{remark}
 
Lo que importa entonces es la \textbf{diferencia} $r_1-r_2$ de las raíces. Esto da tres casos (Figura~\ref{fig:s4-raices}):
 
\begin{enumerate}[label=\textbf{\Roman*.}, leftmargin=2.8em]
    \item $r_1-r_2$ \emph{no es un entero} (incluye raíces reales distintas que no difieren en un entero y raíces complejas conjugadas).
    \item $r_1=r_2$ (raíz doble).
    \item $r_1-r_2=N$, con $N$ un entero positivo.
\end{enumerate}
 
\begin{figure}[H]
\centering
\begin{tikzpicture}[x=1.15cm, lab/.style={anchor=east, text=black, font=\small}]
  \node[lab] at (-.4,2.0) {\textbf{Caso I}};
  \draw[bbegray, ->] (-.2,2.0) -- (6.4,2.0);
  \fill[bbeblue] (1.2,2.0) circle (2.6pt) node[above=4pt, text=black, font=\footnotesize] {$r_2$};
  \fill[bbeblue] (3.9,2.0) circle (2.6pt) node[above=4pt, text=black, font=\footnotesize] {$r_1$};
  \draw[<->, bbeamber, thick] (1.2,1.7) -- (3.9,1.7) node[midway, below, text=black, font=\footnotesize] {$r_1-r_2\notin\mathbb Z$};
  \node[lab] at (-.4,0.6) {\textbf{Caso II}};
  \draw[bbegray, ->] (-.2,0.6) -- (6.4,0.6);
  \fill[bbeblue] (2.8,0.6) circle (2.6pt) node[above=4pt, text=black, font=\footnotesize] {$r_1=r_2$};
  \node[lab] at (-.4,-0.8) {\textbf{Caso III}};
  \draw[bbegray, ->] (-.2,-0.8) -- (6.4,-0.8);
  \foreach \t in {0,1,...,5} \draw[bbegray] (\t+.4,-.88) -- (\t+.4,-.72);
  \fill[bbeblue] (1.4,-0.8) circle (2.6pt) node[above=4pt, text=black, font=\footnotesize] {$r_2$};
  \fill[bbeblue] (3.4,-0.8) circle (2.6pt) node[above=4pt, text=black, font=\footnotesize] {$r_1=r_2+N$};
  \draw[<->, bbered, thick] (1.4,-1.1) -- (3.4,-1.1) node[midway, below, text=black, font=\footnotesize] {$N=2$ (ejemplo)};
\end{tikzpicture}
\caption{Posición relativa de las raíces de la ecuación indicial en cada caso (en el Caso III, $F(r_2+N)=F(r_1)=0$).}
\label{fig:s4-raices}
\end{figure}
 
Enunciaremos el resultado general en tres partes, una para cada caso; recuerden que el problema que estamos tratando es el de encontrar una segunda solución linealmente independiente.
 
\begin{theorem}[Caso I: $r_1-r_2\notin\mathbb Z$]\label{thm:s4-I}
Establecer $r = r_2$ en \eqref{eq:s4-recgen} produce
\[ F(r_2 + n)a_n = -\sum_{k=0}^{n-1} [p_{n-k}(k + r_2) + q_{n-k}] a_k. \]
Como los únicos ceros de $F$ son $r_1$ y $r_2$, se tiene $F(r_2+n)=0$ si y solo si $r_2+n=r_1$, lo cual es imposible si $r_1-r_2$ no es entero. Por lo tanto existe una segunda solución en serie de Frobenius
\[ y_2(x) = x^{r_2} \left[ 1 + \sum_{n=1}^{\infty} a_n(r_2)x^n \right], \]
donde $a_n(r_2)$ denota los valores de los coeficientes obtenidos de \eqref{eq:s4-recgen} cuando $r = r_2$, con $a_0=1$. Las soluciones $y_1$ y $y_2$ son linealmente independientes en (al menos) $0 < x < R$. Si las raíces son complejas conjugadas, $y_1$ es de valor complejo y, por linealidad, sus partes real e imaginaria son soluciones reales linealmente independientes.
\end{theorem}
 
\begin{theorem}[Caso II: $r_1=r_2$]\label{thm:s4-II}
En este caso, ciertamente \textit{no puede} existir una segunda solución en \textit{serie de Frobenius} linealmente independiente. Sin embargo, la técnica de reducción de orden se puede usar para establecer la existencia de una segunda solución linealmente independiente de la forma
\[ y_2(x) = y_1(x) \ln x + x^{r_1} \sum_{n=1}^{\infty} b_n x^n \]
válida para (al menos) $0 < x < R$. Los coeficientes $b_n$ pueden obtenerse sustituyendo esta expresión para $y_2$ en la ecuación diferencial \eqref{eq:s4-frobenius}.
\end{theorem}
 
\begin{theorem}[Caso III: $r_1-r_2=N\in\mathbb N$]\label{thm:s4-III}
Existe una solución de Frobenius $y_1$ para $r_1$, y una segunda solución linealmente independiente de la forma
\[ y_2(x) = A y_1(x) \ln x + x^{r_2} \sum_{n=0}^{\infty} b_n x^n, \qquad b_0\ne0, \]
válida en $(0,R)$, donde las constantes $A$ y $b_n$ se determinan por sustitución directa. La constante $A$ puede ser $0$ o no; lo decide la condición de consistencia \eqref{eq:s4-consistencia} (ver abajo).
\end{theorem}
 
\textbf{¿De dónde sale el logaritmo en el Caso III?} Sustituyendo $r_2=r_1-N$ en \eqref{eq:s4-recgen} y haciendo $n=N$, el coeficiente de $a_N$ vale $F(r_1)=0$ y queda la \emph{condición de consistencia}
\begin{equation}
0 \cdot a_N = -\sum_{k=0}^{N-1} [p_{N-k}(k + r_2) + q_{N-k}] a_k.
\label{eq:s4-consistencia}
\end{equation}
Dado que todos los coeficientes $a_1, a_2, \dots, a_{N-1}$ ya se habrán determinado en términos de $a_0$, \eqref{eq:s4-consistencia} será de la forma
\begin{equation}
0 \cdot a_N = \alpha a_0,
\label{eq:s4-alfa}
\end{equation}
donde $\alpha$ es una constante. Por suposición, $a_0$ es distinto de cero, y por lo tanto surgen dos posibilidades. (Note que \eqref{eq:s4-alfa} \textit{no} es una ecuación para determinar $\alpha$; más bien, es una condición de consistencia para la validez de la relación de recurrencia.)
 
\begin{enumerate}
    \item[\textbf{(a)}] Primero puede ocurrir que $\alpha = 0$. Si esto ocurre, entonces $a_N$ puede especificarse arbitrariamente, y \eqref{eq:s4-recgen} determina todos los coeficientes de Frobenius restantes. Por lo tanto, \textit{sí} obtenemos una segunda solución en serie de Frobenius linealmente independiente (y $A=0$).
    \item[\textbf{(b)}] En el caso más general, $\alpha$ será distinto de cero. Entonces, dado que $a_0 \neq 0$, \eqref{eq:s4-alfa} \textit{no se puede} satisfacer, y por lo tanto no podemos calcular los coeficientes de Frobenius: \textit{no existe} una solución en serie de Frobenius correspondiente a $r = r_2$. La reducción de orden prueba entonces que $y_2$ tiene la forma con $\ln x$ y $A\neq0$. La derivación es directa, pero bastante larga.
\end{enumerate}
La forma del Teorema~\ref{thm:s4-III} incluye el caso (a), que surge cuando $A = 0$.
 
\begin{figure}[H]
\centering
\begin{tikzpicture}[node distance=.7cm and 1.1cm,
   dec/.style={bbecaja, fill=bbeamber!18, draw=bbeamber, text width=4.6cm},
   res/.style={bbecaja, text width=3.6cm, font=\small\sffamily}]
  \node[bbecaja, text width=4.6cm] (a) {Calcular las raíces de $F(r)=0$};
  \node[dec, below=of a] (b) {¿$r_1-r_2$ es un entero?};
  \node[res, right=of b, fill=bbegreen!15, draw=bbegreen] (c1) {\textbf{Caso I}: dos series de Frobenius};
  \node[dec, below=of b] (c) {¿$r_1=r_2$?};
  \node[res, right=of c, fill=bbeblue!15, draw=bbeblue] (c2) {\textbf{Caso II}: $y_2=y_1\ln x+x^{r}\sum_{n\ge1}b_nx^n$};
  \node[dec, below=of c] (d) {\textbf{Caso III}: calcular $\alpha$ en \eqref{eq:s4-alfa}. ¿$\alpha=0$?};
  \node[res, right=of d, fill=bbegreen!15, draw=bbegreen] (c3) {$A=0$: segunda serie de Frobenius};
  \node[res, below=of d, fill=bbered!12, draw=bbered] (c4) {$A\neq0$: $y_2=Ay_1\ln x+x^{r_2}\sum b_nx^n$};
  \draw[bbeflujo] (a) -- (b);
  \draw[bbeflujo] (b) -- node[above, text=black, font=\footnotesize] {no} (c1);
  \draw[bbeflujo] (b) -- node[left, text=black, font=\footnotesize] {sí} (c);
  \draw[bbeflujo] (c) -- node[above, text=black, font=\footnotesize] {sí} (c2);
  \draw[bbeflujo] (c) -- node[left, text=black, font=\footnotesize] {no} (d);
  \draw[bbeflujo] (d) -- node[above, text=black, font=\footnotesize] {sí} (c3);
  \draw[bbeflujo] (d) -- node[left, text=black, font=\footnotesize] {no} (c4);
\end{tikzpicture}
\caption{Diagrama de decisión para la forma de la segunda solución.}
\label{fig:s4-decision}
\end{figure}
 
A modo de resumen:
 
\begin{remark}
  Considere la ecuación diferencial
\begin{equation}
x^2 y'' + xp(x)y' + q(x)y = 0,
\label{eq:s4-resumen}
\end{equation}
donde $p$ y $q$ son analíticas en $x = 0$. Suponga que
\[ p(x) = \sum_{n=0}^{\infty} p_n x^n, \qquad q(x) = \sum_{n=0}^{\infty} q_n x^n, \]
para $|x| < R$. Sean $r_1$ y $r_2$ las raíces de la ecuación indicial y asuma que $r_1 \ge r_2$ si estas raíces son reales. Entonces \eqref{eq:s4-resumen} tiene dos soluciones linealmente independientes válidas (al menos) en el intervalo $(0, R)$. La forma de la solución se determina de la siguiente manera:
 
\begin{enumerate}
    \item $r_1 - r_2$ no es un entero:
    \begin{align}
    y_1(x) &= x^{r_1} \sum_{n=0}^{\infty} a_n x^n, \qquad a_0 \neq 0, \label{eq:s4-I-1}\\
    y_2(x) &= x^{r_2} \sum_{n=0}^{\infty} b_n x^n, \qquad b_0 \neq 0. \label{eq:s4-I-2}
    \end{align}
    \item $r_1 = r_2 = r$:
    \begin{align}
    y_1(x) &= x^r \sum_{n=0}^{\infty} a_n x^n, \qquad a_0 \neq 0, \label{eq:s4-II-1}\\
    y_2(x) &= y_1(x) \ln x + x^r \sum_{n=1}^{\infty} b_n x^n. \label{eq:s4-II-2}
    \end{align}
    \item $r_1 - r_2$ es un entero positivo:
    \begin{align}
    y_1(x) &= x^{r_1} \sum_{n=0}^{\infty} a_n x^n, \qquad a_0 \neq 0, \label{eq:s4-III-1}\\
    y_2(x) &= A y_1(x) \ln x + x^{r_2} \sum_{n=0}^{\infty} b_n x^n, \qquad b_0 \neq 0. \label{eq:s4-III-2}
    \end{align}
\end{enumerate}
 
Los coeficientes en cada una de estas soluciones pueden determinarse por sustitución directa en la Ecuación \eqref{eq:s4-resumen}. Finalmente, si $x^{r_1}$ y $x^{r_2}$ se reemplazan por $|x|^{r_1}$ y $|x|^{r_2}$, respectivamente, obtenemos soluciones linealmente independientes que son válidas para (al menos) $0 < |x| < R$.
\end{remark}
 
\begin{table}[H]
\centering
\small
\caption{Resumen de los tres casos del método de Frobenius.}
\label{tab:s4-casos}
\begin{tabular}{|c|c|p{3.2cm}|p{4.6cm}|p{2.8cm}|}
\hline
\rowcolor{bbe!15}\textbf{Caso} & \textbf{Condición} & $\mathbf{y_1}$ & $\mathbf{y_2}$ & \textbf{¿Aparece $\ln x$?} \\ \hline
I & $r_1-r_2\notin\mathbb Z$ & $x^{r_1}\sum a_nx^n$ & $x^{r_2}\sum b_nx^n$ & no \\ \hline
II & $r_1=r_2=r$ & $x^{r}\sum a_nx^n$ & $y_1\ln x+x^{r}\sum_{n\ge1}b_nx^n$ & siempre \\ \hline
III & $r_1-r_2=N\in\mathbb N$ & $x^{r_1}\sum a_nx^n$ & $Ay_1\ln x+x^{r_2}\sum b_nx^n$ & solo si $\alpha\neq0$ ($A\ne0$) \\ \hline
\end{tabular}
\end{table}
 
Revisemos algunos ejemplos para ilustrar lo anterior:
 
\begin{example}\label{ex:s4-doble}
    Considere la ecuación diferencial
\begin{equation}
x^2 y'' - x(3 + x)y' + (4 - x)y = 0.
\label{eq:s4-ej1}
\end{equation}
Determine la forma general de dos soluciones en serie linealmente independientes en la vecindad del punto singular regular $x = 0$.
 
\noindent \textbf{Solución:} \quad Por inspección vemos que $x = 0$ es un punto singular regular de la Ecuación \eqref{eq:s4-ej1} y que la ecuación indicial es
\[ r(r - 1) - 3r + 4 = r^2 - 4r + 4 = (r - 2)^2 = 0. \]
Dado que $r = 2$ es una raíz repetida (Caso II), se deduce del Teorema~\ref{thm:s4-II} que existen dos soluciones linealmente independientes para la Ecuación \eqref{eq:s4-ej1} de la forma
\[ y_1(x) = x^2 \sum_{n=0}^{\infty} a_n x^n, \qquad y_2(x) = y_1(x) \ln |x| + x^2 \sum_{n=1}^{\infty} b_n x^n. \]
 
En este problema, tenemos
\[ p(x) = -(3 + x), \qquad q(x) = 4 - x. \]
Dado que estos son polinomios, sus expansiones en serie de potencias alrededor de $x = 0$ convergen para todo $x$ real. Se deduce que las soluciones en serie de la Ecuación \eqref{eq:s4-ej1} serán válidas para todo $x \neq 0$.
 
\vspace{.2cm}
\noindent\textit{Los coeficientes.} Con $p_1=-1$, $q_1=-1$ y $F(r+n)=(r+n-2)^2$, la fórmula \eqref{eq:s4-recgen} tiene una sola contribución ($k=n-1$):
\[ (r+n-2)^2a_n=-\bigl[(r+n-1)p_1+q_1\bigr]a_{n-1}=(r+n)\,a_{n-1}. \]
Con $r=2$ resulta $n^2a_n=(n+2)a_{n-1}$ y por lo tanto $a_n=\dfrac{(n+2)!}{2\,(n!)^2}$ (con $a_0=1$), es decir
\[ y_1(x)=x^2\Bigl(1+3x+3x^2+\tfrac53x^3+\tfrac{5}{8}x^4+\cdots\Bigr). \]
\hfill $\square$
\end{example}
 
\begin{remark}
\textbf{Un truco para el Caso II.} Si $a_n(r)$ son los coeficientes de \eqref{eq:s4-recgen} \emph{como funciones de $r$} (con $a_0=1$), los coeficientes del logaritmo se obtienen derivando respecto de $r$:
\[ b_n=\frac{d\,a_n}{dr}\Big|_{r=r_1}. \]
En el Ejemplo~\ref{ex:s4-doble}, $a_n(r)=\prod_{k=1}^{n}\dfrac{r+k}{(r+k-2)^2}$, y la derivada logarítmica da
\[ b_n=a_n(2)\sum_{k=1}^{n}\Bigl(\frac1{k+2}-\frac2k\Bigr),\qquad b_1=-5,\ \ b_2=-\tfrac{29}{4},\ \ b_3=-\tfrac{173}{36}. \]
\end{remark}
 
\begin{figure}[H]
\centering
\begin{tikzpicture}
\begin{axis}[bbeserie, height=6cm, xmin=0, xmax=1.2, ymin=-22, ymax=22, restrict y to domain=-40:40,
  legend pos=north west, samples=100]
\addplot[bbeblue, very thick, smooth] coordinates {(0.02,0.0004) (0.04,0.0018) (0.06,0.0043) (0.08,0.0081) (0.1,0.0133) (0.12,0.0202) (0.14,0.0291) (0.16,0.0400) (0.18,0.0534) (0.2,0.0694) (0.22,0.0883) (0.24,0.1105) (0.26,0.1362) (0.28,0.1659) (0.3,0.1998) (0.32,0.2385) (0.34,0.2822) (0.36,0.3315) (0.38,0.3869) (0.4,0.4487) (0.42,0.5177) (0.44,0.5942) (0.46,0.6790) (0.48,0.7727) (0.5,0.8759) (0.52,0.9893) (0.54,1.1138) (0.56,1.2500) (0.58,1.3988) (0.6,1.5612) (0.62,1.7380) (0.64,1.9302) (0.66,2.1388) (0.68,2.3650) (0.7,2.6099) (0.72,2.8747) (0.74,3.1606) (0.76,3.4691) (0.78,3.8014) (0.8,4.1591) (0.82,4.5437) (0.84,4.9569) (0.86,5.4003) (0.88,5.8758) (0.9,6.3853) (0.92,6.9306) (0.94,7.5139) (0.96,8.1374) (0.98,8.8033) (1,9.5140) (1.02,10.2720) (1.04,11.0799) (1.06,11.9405) (1.08,12.8567) (1.1,13.8313) (1.12,14.8677) (1.14,15.9689) (1.16,17.1386) (1.18,18.3802) (1.2,19.6976)};
\addplot[bbered, very thick, smooth] coordinates {(0.02,-0.0017) (0.04,-0.0061) (0.06,-0.0132) (0.08,-0.0232) (0.1,-0.0364) (0.12,-0.0532) (0.14,-0.0740) (0.16,-0.0991) (0.18,-0.1293) (0.2,-0.1649) (0.22,-0.2067) (0.24,-0.2551) (0.26,-0.3109) (0.28,-0.3748) (0.3,-0.4476) (0.32,-0.5301) (0.34,-0.6232) (0.36,-0.7277) (0.38,-0.8448) (0.4,-0.9753) (0.42,-1.1205) (0.44,-1.2815) (0.46,-1.4596) (0.48,-1.6560) (0.5,-1.8722) (0.52,-2.1096) (0.54,-2.3699) (0.56,-2.6546) (0.58,-2.9654) (0.6,-3.3043) (0.62,-3.6731) (0.64,-4.0740) (0.66,-4.5089) (0.68,-4.9803) (0.7,-5.4904) (0.72,-6.0418) (0.74,-6.6371) (0.76,-7.2792) (0.78,-7.9708) (0.8,-8.7151) (0.82,-9.5153) (0.84,-10.3748) (0.86,-11.2972) (0.88,-12.2860) (0.9,-13.3454) (0.92,-14.4793) (0.94,-15.6921) (0.96,-16.9882) (0.98,-18.3725) (1,-19.8498) (1.02,-21.4254) (1.04,-23.1046) (1.06,-24.8932) (1.08,-26.7972) (1.1,-28.8227) (1.12,-30.9762) (1.14,-33.2647) (1.16,-35.6952) (1.18,-38.2751)};
\legend{$y_1$, $y_2=y_1\ln x+x^2\sum b_nx^n$}
\end{axis}
\end{tikzpicture}
\caption{Ejemplo~\ref{ex:s4-doble} (raíz doble $r=2$): ambas soluciones tienden a $0$ en $x=0$, pero $y_2\sim x^2\ln x$ lo hace con una derivada segunda no acotada; el logaritmo es la huella del Caso II.}
\label{fig:s4-doble}
\end{figure}
 
\begin{example}\label{ex:s4-gamma}
    Considere la ecuación diferencial
\begin{equation}
x^2 y'' + x(1 + 2x)y' - \frac{1}{4}(1 - \gamma x)y = 0, \qquad x > 0,
\label{eq:s4-ej2}
\end{equation}
donde $\gamma$ es una constante. Determine la forma de dos soluciones en serie linealmente independientes para esta ecuación diferencial alrededor del punto singular regular $x = 0$.
 
\noindent \textbf{Solución:} \quad En este caso, $p(x)=1+2x$ y $q(x)=-\frac14+\frac{\gamma}{4}x$; luego $p_0 = 1$, $p_1=2$, $q_0 = -\frac{1}{4}$ y $q_1=\frac\gamma4$, de modo que
\[ F(r) = r(r - 1) + r - \frac{1}{4} = r^2 - \frac{1}{4}. \]
Así, las raíces de la ecuación indicial son $r = \pm \frac{1}{2}$. Estableciendo $r_1 = \frac{1}{2}$ y $r_2 = -\frac{1}{2}$, se deduce que $r_1 - r_2 = 1$, de modo que estamos en el Caso III. Así, existen dos soluciones linealmente independientes para la Ecuación \eqref{eq:s4-ej2} en $(0, \infty)$ de la forma
\begin{equation}
y_1(x) = x^{1/2} \sum_{n=0}^{\infty} a_n x^n, \qquad y_2(x) = A y_1(x) \ln x + x^{-1/2} \sum_{n=0}^{\infty} b_n x^n.
\label{eq:s4-formas-gamma}
\end{equation}
Para determinar si la constante $A$ es cero o distinta de cero, necesitamos la relación de recurrencia general. Sustituir $y=x^r\sum a_nx^n$ en \eqref{eq:s4-ej2} produce (aquí \eqref{eq:s4-recgen} tiene solo el término $k=n-1$)
\[ \bigl[(r+n)^2-\tfrac14\bigr]a_n=-\Bigl[2(r+n-1)+\frac{\gamma}{4}\Bigr]a_{n-1},\qquad n=1,2,\dots \]
Cuando $r=-\frac12$, esto se reduce a
\begin{equation}
n(n-1)a_n=-\Bigl(2n-3+\frac{\gamma}{4}\Bigr)a_{n-1}.
\label{eq:s4-rec-gamma}
\end{equation}
Como predice nuestra teoría general, el coeficiente de $a_n$ es cero cuando $n=r_1-r_2=1$. Así, existe una segunda solución en serie de Frobenius [$A=0$ en \eqref{eq:s4-formas-gamma}] si y solo si el término del lado derecho también se anula cuando $n=1$. Establecer $n=1$ en \eqref{eq:s4-rec-gamma} produce la condición de consistencia
\begin{equation}
0\cdot a_1=\Bigl(1-\frac{\gamma}{4}\Bigr)a_0.
\label{eq:s4-consistencia-gamma}
\end{equation}
Dado que $a_0\neq0$, se deduce de \eqref{eq:s4-consistencia-gamma} que cuando $\gamma\neq4$, no existe una segunda solución en serie de Frobenius linealmente independiente, y por lo tanto la constante $A$ en \eqref{eq:s4-formas-gamma} es necesariamente distinta de cero. Si $\gamma=4$, sin embargo, entonces \eqref{eq:s4-consistencia-gamma} se satisface idénticamente independientemente del valor de $a_1$. En este caso, podemos especificar $a_1$ arbitrariamente y entonces la relación de recurrencia \eqref{eq:s4-rec-gamma} determinará los coeficientes restantes en una segunda solución en serie de Frobenius linealmente independiente. En resumen:
 
\begin{enumerate}
    \item Si $\gamma\neq4$, existen dos soluciones linealmente independientes de la forma \eqref{eq:s4-formas-gamma} con $A$ necesariamente distinto de cero (con $a_0=b_0=1$ resulta $A=\frac{4-\gamma}{4}$; Tabla~\ref{tab:s4-A}).
    \item Si $\gamma=4$, entonces existen dos soluciones en serie de Frobenius linealmente independientes
    \[ y_1(x) = x^{1/2} \sum_{n=0}^{\infty} a_n x^n, \qquad y_2(x) = x^{-1/2} \sum_{n=0}^{\infty} b_n x^n. \]
\end{enumerate}
En ambos casos, las soluciones en serie serán válidas para $0 < x < \infty$. \hfill $\square$
\end{example}
 
\begin{table}[H]
\centering
\small
\caption{Ejemplo~\ref{ex:s4-gamma}: la constante $A$ del logaritmo según $\gamma$ (con $a_0=b_0=1$).}
\label{tab:s4-A}
\begin{tabular}{|c|c|c|c|c|c|}
\hline
\rowcolor{bbe!15}$\gamma$ & $0$ & $1$ & $2$ & $4$ & $8$ \\ \hline
$A=\dfrac{4-\gamma}{4}$ & $1$ & $\frac34$ & $\frac12$ & $\mathbf{0}$ & $-1$ \\ \hline
¿Hay $\ln x$? & sí & sí & sí & \textbf{no} & sí \\ \hline
\end{tabular}
\end{table}
 
\begin{remark}
\textbf{El caso $\gamma=4$ se puede resolver a mano.} La ecuación queda $x^2y''+x(1+2x)y'+\bigl(x-\tfrac14\bigr)y=0$ y se verifica directamente que $x^{-1/2}$ y $x^{-1/2}e^{-2x}$ son soluciones. Por lo tanto la solución general es $y=x^{-1/2}\bigl(c_1+c_2e^{-2x}\bigr)$. La serie de Frobenius para $r_1=\tfrac12$ es $y_1=(1-e^{-2x})/(2\sqrt x)$, y para $r_2=-\tfrac12$ es $y_2=e^{-2x}/\sqrt x=x^{-1/2}\sum_{n}\frac{(-2)^n}{n!}x^n$. La Figura~\ref{fig:s4-gamma4} las muestra.
\end{remark}
 
\begin{figure}[H]
\centering
\begin{tikzpicture}
\begin{axis}[bbeserie, height=5.6cm, xmin=0, xmax=3, ymin=0, ymax=3, restrict y to domain=-1:8,
  legend pos=north east, samples=120]
\addplot[bbeblue, very thick, domain=0.03:3] {(1-exp(-2*x))/(2*sqrt(x))};
\addplot[bbered, very thick, domain=0.03:3] {exp(-2*x)/sqrt(x)};
\legend{$y_1=(1-e^{-2x})/(2\sqrt x)$, $y_2=e^{-2x}/\sqrt x$}
\end{axis}
\end{tikzpicture}
\caption{Ejemplo~\ref{ex:s4-gamma} con $\gamma=4$ ($A=0$): dos series de Frobenius, con exponentes $\frac12$ y $-\frac12$.}
\label{fig:s4-gamma4}
\end{figure}
 
\begin{example}\label{ex:s4-log}
    Determine dos soluciones en serie linealmente independientes para
\begin{equation}
x^2 y'' + x(3 - x)y' + y = 0, \qquad x > 0.
\label{eq:s4-ej3}
\end{equation}
 
\noindent \textbf{Solución:} \quad Vemos por inspección que $x = 0$ es un punto singular regular. Además, dado que las funciones
\[ p(x) = 3 - x \quad \text{y} \quad q(x) = 1 \]
son polinomios, las soluciones en serie linealmente independientes obtenidas en el Teorema~\ref{thm:s4-II} serán válidas en $(0, \infty)$. Comenzamos por obtener una solución en serie de Frobenius. Sustituir
\[ y(x) = x^r \sum_{n=0}^{\infty} a_n x^n, \qquad a_0 \neq 0, \]
en \eqref{eq:s4-ej3} produce
\begin{align*}
\sum_{n=0}^{\infty} (r + n)(r + n - 1)a_n x^{r+n} &+ 3\sum_{n=0}^{\infty} (r + n)a_n x^{r+n} \\
&- \sum_{n=0}^{\infty} (r + n)a_n x^{r+n+1} + \sum_{n=0}^{\infty} a_n x^{r+n} = 0,
\end{align*}
lo cual, al recolectar los coeficientes de $x^{r+n}$ y reemplazar $n$ por $n - 1$ en la tercera suma, y dividir por $x^r$, se puede escribir como
\begin{equation}
\sum_{n=0}^{\infty} [(r + n)(r + n + 2) + 1]a_n x^n - \sum_{n=1}^{\infty} (r + n - 1)a_{n-1} x^n = 0.
\label{eq:s4-suma-3}
\end{equation}
Determinamos los $a_n$ de la manera habitual. Cuando $n = 0$, obtenemos $[r(r + 2) + 1]a_0 = 0$, de modo que la ecuación indicial es $(r + 1)^2 = 0$. Por lo tanto hay solo una raíz: $r = -1$ (Caso II). De \eqref{eq:s4-suma-3}, los coeficientes restantes deben satisfacer
\[ [(r + n)(r + n + 2) + 1]a_n - (r + n - 1)a_{n-1} = 0, \qquad n = 1, 2, \dots \]
Estableciendo $r = -1$, esto se simplifica a
\[ a_n = \frac{n - 2}{n^2}a_{n-1}, \qquad n = 1, 2, 3, \dots \]
En consecuencia,
\[ a_1 = -a_0, \qquad a_2 = a_3 = a_4 = \dots = 0 \]
(el factor $n-2$ \emph{corta} la serie en $n=2$). Así, una solución en serie de Frobenius para \eqref{eq:s4-ej3} es
\begin{equation}
y_1(x) = x^{-1}(1 - x),
\label{eq:s4-y1}
\end{equation}
donde establecemos $a_0 = 1$.
 
Dado que la ecuación indicial tiene solo una raíz, no existe una segunda solución en serie de Frobenius linealmente independiente. Sin embargo, de acuerdo con el Teorema~\ref{thm:s4-II}, existe una segunda solución linealmente independiente de la forma
\begin{equation}
y_2(x) = y_1(x) \ln x + x^{-1} \sum_{n=1}^{\infty} b_n x^n,
\label{eq:s4-y2}
\end{equation}
donde los coeficientes $b_n$ se pueden determinar por sustitución en \eqref{eq:s4-ej3}. Los cálculos son bastante directos, pero son largos y tediosos. Se anima a prestar toda su atención y a no dejarse abrumar por el aspecto de las ecuaciones. Derivando \eqref{eq:s4-y2},
\begin{align*}
y_2' &= y_1' \ln x + x^{-1} y_1 + \sum_{n=1}^{\infty} (n - 1)b_n x^{n-2}, \\
y_2'' &= y_1'' \ln x + 2x^{-1}y_1' - x^{-2}y_1 + \sum_{n=1}^{\infty} (n - 1)(n-2)b_n x^{n-3}.
\end{align*}
Sustituyendo en \eqref{eq:s4-ej3}, el factor que multiplica a $\ln x$ es $x^2y_1''+x(3-x)y_1'+y_1=0$ (pues $y_1$ es solución) y se anula. Agrupando lo que queda,
\[ 2xy_1' + 2y_1 - xy_1 + \sum_{n=1}^{\infty} [(n-1)(n-2)+3(n-1)+1]\,b_n x^{n-1} - \sum_{n=1}^{\infty} (n - 1)b_n x^n = 0 . \]
Como $(n-1)(n-2)+3(n-1)+1=n^2$, y reemplazando $n$ por $n-1$ en la última suma,
\begin{equation}
2xy_1' + 2y_1 - xy_1 + \sum_{n=1}^{\infty} n^2 b_n x^{n-1} - \sum_{n=2}^{\infty} (n - 2)b_{n-1}x^{n-1} = 0.
\label{eq:s4-28}
\end{equation}
De \eqref{eq:s4-y1}, tenemos $y_1(x) = x^{-1}(1 - x)$ e $y_1'(x) = -x^{-2}$. Sustituyendo estas expresiones en \eqref{eq:s4-28} da
\[ -2x^{-1} + 2x^{-1}(1 - x) - (1 - x) + \sum_{n=1}^{\infty} n^2 b_n x^{n-1} - \sum_{n=2}^{\infty} (n - 2)b_{n-1}x^{n-1} = 0, \]
es decir,
\[ -3 + x + \sum_{n=1}^{\infty} n^2 b_n x^{n-1} - \sum_{n=2}^{\infty} (n - 2)b_{n-1}x^{n-1} = 0. \]
Igualando los coeficientes correspondientes de $x^{n-1}$ a cero para $n = 1, 2, \dots$, obtenemos los $b_n$ de manera familiar:
\begin{align*}
n &= 1: & -3 + b_1 &= 0, \text{ lo que implica que } b_1 = 3, \\
n &= 2: & 1 + 4b_2 &= 0, \text{ lo que implica que } b_2 = -\frac{1}{4}, \\
n &\ge 3: & n^2 b_n &- (n - 2)b_{n-1} = 0.
\end{align*}
Es decir,
\begin{equation}
b_n = \frac{n - 2}{n^2}b_{n-1}, \qquad n = 3, 4, \dots
\label{eq:s4-29}
\end{equation}
Cuando $n = 3$, $b_3 = \frac{1}{3^2}b_2 = -\frac{1}{2^2 \cdot 3^2}$. Cuando $n = 4$, \eqref{eq:s4-29} implica que $b_4 = \frac{2}{4^2}b_3 = -\frac{1 \cdot 2}{2^2 \cdot 3^2 \cdot 4^2}$. Vemos que en general, para $n \ge 2$,
\[ b_n = -\frac{(n - 2)!}{2^2 \cdot 3^2 \cdots n^2} = -\frac{(n - 2)!}{(n!)^2} \]
(la fórmula también vale para $n=2$: $-\frac{0!}{(2!)^2}=-\frac14$). Finalmente, sustituir los $b_n$ en \eqref{eq:s4-y2} produce la siguiente solución a la Ecuación \eqref{eq:s4-ej3}:
\[ y_2(x) = x^{-1}(1 - x) \ln x + 3 - \sum_{n=2}^{\infty} \frac{(n - 2)!}{(n!)^2}\, x^{n-1}. \]
El Teorema~\ref{thm:s4-II} garantiza que $y_1$ y $y_2$ son linealmente independientes en $(0, \infty)$. \hfill $\square$
\end{example}
 
\begin{figure}[H]
\centering
\begin{tikzpicture}
\begin{axis}[bbeserie, height=6cm, xmin=0, xmax=3, ymin=-4, ymax=6, restrict y to domain=-12:12,
  legend pos=north east, samples=150]
\addplot[bbeblue, very thick, domain=0.15:3] {(1-x)/x};
\addplot[bbered, very thick, smooth] coordinates {(0.15,-7.7885) (0.2,-3.4889) (0.25,-1.2232) (0.3,0.1131) (0.35,0.9593) (0.4,1.5209) (0.45,1.9056) (0.5,2.1744) (0.55,2.3643) (0.6,2.4986) (0.65,2.5928) (0.7,2.6572) (0.75,2.6994) (0.8,2.7245) (0.85,2.7364) (0.9,2.7380) (0.95,2.7314) (1,2.7183) (1.05,2.7000) (1.1,2.6774) (1.15,2.6514) (1.2,2.6226) (1.25,2.5915) (1.3,2.5585) (1.35,2.5239) (1.4,2.4880) (1.45,2.4510) (1.5,2.4131) (1.55,2.3744) (1.6,2.3351) (1.65,2.2952) (1.7,2.2549) (1.75,2.2142) (1.8,2.1731) (1.85,2.1317) (1.9,2.0900) (1.95,2.0481) (2,2.0060) (2.05,1.9637) (2.1,1.9212) (2.15,1.8785) (2.2,1.8357) (2.25,1.7927) (2.3,1.7495) (2.35,1.7062) (2.4,1.6626) (2.45,1.6189) (2.5,1.5751) (2.55,1.5310) (2.6,1.4867) (2.65,1.4423) (2.7,1.3976) (2.75,1.3526) (2.8,1.3075) (2.85,1.2620) (2.9,1.2163) (2.95,1.1704) (3,1.1241)};
\legend{$y_1=(1-x)/x$, $y_2=y_1\ln x+\cdots$}
\end{axis}
\end{tikzpicture}
\caption{Ejemplo~\ref{ex:s4-log}: $y_1\sim1/x\to+\infty$ y $y_2\sim(\ln x)/x\to-\infty$ cuando $x\to0^+$; el logaritmo cambia el signo y acelera la explosión.}
\label{fig:s4-log}
\end{figure}
 
El Caso III no siempre trae un logaritmo, aunque las raíces difieran en un entero. El ejemplo siguiente (que reaparecerá como la ecuación de Bessel de orden $\tfrac12$) muestra el caso $\alpha=0$.
 
\begin{example}\label{ex:s4-sin-log}
\textbf{(Caso III sin logaritmo.)} Determine dos soluciones de $x^2y''+xy'+\bigl(x^2-\tfrac14\bigr)y=0$, $x>0$.
 
\noindent \textbf{Solución:} \quad Aquí $p\equiv1$ y $q(x)=x^2-\frac14$ (así $p_0=1$, $q_0=-\tfrac14$, $q_2=1$ y todos los demás $p_n,q_n$ son $0$). La ecuación indicial es $F(r)=r^2-\tfrac14=0$, con $r_1=\tfrac12$ y $r_2=-\tfrac12$ (diferencia $N=1$: Caso III). Condición de consistencia \eqref{eq:s4-consistencia} con $N=1$:
\[ 0\cdot a_1=-\bigl[(r_2+0)\,p_1+q_1\bigr]a_0=0, \]
pues $p_1=q_1=0$. Es decir, $\alpha=0$: \emph{no hay logaritmo} y $a_1$ queda libre. Como en \eqref{eq:s4-recgen} solo sobrevive $q_2=1$, la recurrencia es $F(r+n)a_n=-a_{n-2}$.
\begin{itemize}
\item $r=\tfrac12$: $F(\tfrac12+n)=n(n+1)$, luego $a_1=0$ y $a_{2m}=\dfrac{(-1)^m}{(2m+1)!}a_0$. Resulta $y_1=x^{1/2}\sum_{m\ge0}\frac{(-1)^m x^{2m}}{(2m+1)!}=\dfrac{\sin x}{\sqrt x}$.
\item $r=-\tfrac12$: $F(-\tfrac12+n)=n(n-1)$, luego $a_{2m}=\dfrac{(-1)^m}{(2m)!}a_0$ y $a_1$ libre (tomar $a_1=0$; un valor distinto de cero solo suma un múltiplo de $y_1$). Resulta $y_2=x^{-1/2}\sum_{m\ge0}\frac{(-1)^m x^{2m}}{(2m)!}=\dfrac{\cos x}{\sqrt x}$.
\end{itemize}
Por lo tanto, la solución general es $y=x^{-1/2}\bigl(c_1\sin x+c_2\cos x\bigr)$. \hfill $\square$
\end{example}
 
\begin{figure}[H]
\centering
\begin{tikzpicture}
\begin{axis}[bbeosc, height=6cm, xmin=0, xmax=14, ymin=-2, ymax=2.5, restrict y to domain=-6:6,
  legend pos=north east, legend columns=2, samples=300]
\bbezero{0}{14}
\addplot[bbeblue, very thick, domain=0.05:14] {sin(deg(x))/sqrt(x)};
\addplot[bbered, very thick, domain=0.05:14] {cos(deg(x))/sqrt(x)};
\addplot[bbegray, dashed, domain=0.05:14] {1/sqrt(x)};
\addplot[bbegray, dashed, domain=0.05:14, forget plot] {-1/sqrt(x)};
\legend{$\sin x/\sqrt x$, $\cos x/\sqrt x$, $\pm1/\sqrt x$}
\end{axis}
\end{tikzpicture}
\caption{Ejemplo~\ref{ex:s4-sin-log}: con exponentes $\pm\frac12$, una solución es acotada en $0$ y la otra explota como $x^{-1/2}$; ambas oscilan con amplitud $\sim x^{-1/2}$.}
\label{fig:s4-sin-log}
\end{figure}
 
\begin{table}[H]
\centering
\small
\caption{Raíces que difieren en un entero: ¿aparece el logaritmo?}
\label{tab:s4-log}
\begin{tabular}{|p{6.4cm}|c|c|c|}
\hline
\rowcolor{bbe!15}\textbf{Ecuación} & $\mathbf{N}$ & $\boldsymbol{\alpha}$ & $\mathbf{A}$ \\ \hline
$x^2y''+x(1+2x)y'-\frac14(1-\gamma x)y=0$, $\gamma\ne4$ & $1$ & $\ne0$ & $\frac{4-\gamma}{4}\neq0$ \\ \hline
$x^2y''+x(1+2x)y'-\frac14(1-4x)y=0$ & $1$ & $0$ & $0$ \\ \hline
$x^2y''+xy'+(x^2-\frac14)y=0$ (Bessel, $p=\frac12$) & $1$ & $0$ & $0$ \\ \hline
$xy''+y=0$ (Ejercicio) & $1$ & $\ne0$ & $\ne0$ \\ \hline
$x^2y''+xy'+(x^2-1)y=0$ (Bessel, $p=1$) & $2$ & $\ne0$ & $\ne0$ \\ \hline
\end{tabular}
\end{table}
 
\begin{exercise}[Para llevar]
Determina la forma (Caso I, II o III; presencia de $\ln x$) de dos soluciones linealmente independientes:
\begin{enumerate}
\item $xy''+y=0$, $x>0$ (multiplica por $x$ para llegar a la forma de Frobenius).
\item $x^2y''-xy'+(1+x)y=0$, $x>0$.
\end{enumerate}
\emph{Respuestas:} (1) $x^2y''+x\,y=0$, es decir $p\equiv0$, $q(x)=x$: $F(r)=r(r-1)$, $r_1=1$, $r_2=0$, $N=1$. La consistencia es $0\cdot a_1=-q_1a_0=-a_0\neq0$, luego \textbf{hay logaritmo}: $y_1=x\bigl(1-\frac x2+\frac{x^2}{12}-\cdots\bigr)$ e $y_2=-y_1\ln x+1+\cdots$ (con $b_0=1$ resulta $A=-1$). \; (2) $F(r)=r(r-1)-r+1=(r-1)^2$: Caso II, raíz doble $r=1$; $y_1=x\sum a_nx^n$ e $y_2=y_1\ln x+x\sum_{n\ge1}b_nx^n$.
\end{exercise}
